\section{All quasisimple and extraspecial pieces in the entire tensor partner}
\label{sec:quasisimple-partner-support}

\paragraph{Status, provenance, and exact hypothesis.}
HD-FT20--22, reserved question Q2dt, extends HD-FT18--19.
The parent supplied the Duyan--Halasi--Podoski route and found the
stronger Guralnick--Saxl conjugate-generation input.  The precise
complex reduction, perfect-lift argument, and synthesis are written
here by \texttt{/root/g34\_eightfolds}.  D-TD08's extraspecial floor
is the separately authored and independently reviewed result of
\texttt{/root/imprimitive}.  The existing entire-partner tensor
argument is shared prior work.  The author and root have separately read the complete frozen proof, source/status record, and invoked source interfaces; both whole-proof checks pass without repair. No independent origin of the shared inputs is claimed.

The entire partner must have the actual matrix-algebra normalization
specified below.  The result does not assert that an arbitrary
primitive group has already been decomposed into those pieces.
Every local operator is of finite order up to scalar, and
\[
 \nu(T)=\dim W-\max_{z\in\mathbf C^*}\dim\ker(T-zI).
 \tag{FTQ.1}
\]
Actual nonidentity scalars have full moving rank although their
projective support is zero.

\subsection*{HD-FT20. A floor for every represented quasisimple core}

\begin{theorem}[HD-FT20]
Let \(E\subset\operatorname{GL}(W)\) be a finite quasisimple group
acting irreducibly, and set \(a=\dim W\).  Its representation is the
actual faithful image; an abstract covering group's represented
central kernel has already been divided out.  Let \(T\) be a
nonscalar finite-projective-order operator normalizing the projective
image of \(E\).
\begin{enumerate}
\item If \(E/Z(E)\) is not alternating, then \(\nu(T)/a\ge1/80\).
\item If \(E/Z(E)\simeq A_m\), \(m\ge7\), then
\[
                         \frac{\nu(T)}a\ge\frac1{m-1}.
 \tag{FTQ.2}
\]
This includes ordinary, split, and spin representations and every
represented central quotient.
\item For \(E/Z(E)\simeq A_5\) or \(A_6\), one has
\(\nu(T)/a\ge1/46>1/80\).
\end{enumerate}
No tensor indecomposability, original coefficient field, or
essential self-duality assumption enters this local theorem.
\end{theorem}

\begin{proof}
First, projective normalization is enough.  It gives
\(TET^{-1}\subset\mathbf C^*E\).  Since \(E\) is perfect,
\[
 TET^{-1}=[TET^{-1},TET^{-1}]
       \subset[\mathbf C^*E,\mathbf C^*E]=E.
 \tag{FTQ.3}
\]
Applying the same argument to \(T^{-1}\) gives equality.  Thus all
commutators and automorphisms below concern the actual represented
group; no characteristic-kernel assertion about an abstract
universal cover is being used.

For the first assertion, the exact published input is
\cite[Corollary~4.3, p.~142]{BI:DHP}: if \(E_0\) is a coprime
quasisimple irreducible linear group over a finite field, not of
alternating type, and \(Z\) is the scalar group, then
\(\operatorname{MinSupp}(ZE_0)\ge\dim(W)/40\).
This stated scope includes every such represented quasisimple
group; it has no restriction to a generic Schur cover.
We now give the complex-coefficient passage explicitly.

Let \(s\) be the exponent of \(E\).  By
\cite[Section~12.3, Theorem~24 and proof, p.~94]{OG:Serre1977},
the representation has a model over \(K=\mathbf Q(\zeta_s)\).
The split matrix-algebra description in
\cite[Section~12.2, pp.~92--93]{OG:Serre1977} says that the
\(K\)-span of its matrices is \(\operatorname{Mat}_a(K)\).
Choose \(a^2\) group matrices forming a basis of that algebra.
Choose a localization \(\mathcal R\) of the ring of integers of
\(K\), obtained by inverting a nonzero integer, containing all
matrix entries and their inverses as needed, and all \(s\)-th roots.
Enlarge that integer so that \(|E|\), the nonzero basis determinant,
and one nonzero maximal-rank minor for every nonzero matrix
\(g-zI\), \(g\in E\), \(z\in\mu_s\), are units in
\(\mathcal R\).  This is possible by inverting integer norms of
their nonzero integral numerators and their denominators.
There are only finitely many matrices and minors.

Choose a maximal ideal of \(\mathcal R\) of residue characteristic
not dividing that integer.  The residue field \(k\) is finite.
All chosen ranks are unchanged: larger minors vanish identically,
and the chosen rank minor stays nonzero.  In particular the reduced
representation is faithful, and the selected group matrices span
\(\operatorname{Mat}_a(k)\).  It is therefore absolutely
irreducible.  Its image \(E_0\) is isomorphic to \(E\), is
quasisimple and coprime to \(\operatorname{char}k\), and has the
same nonalternating simple quotient.  Applying the cited corollary
to \(k^*E_0\) and to every reduced \(z^{-1}g\ne I\) proves
\[
                 \nu(g)\ge a/40\qquad(g\in E\setminus Z(E)).
 \tag{FTQ.4}
\]
All eigenvalues of \(g\) lie in \(\mu_s\), so the list covers every
target scalar relevant to \(\nu\).  This proves the passage to
characteristic zero, including arbitrary central-cover kernels,
rather than assuming a coprime representation transfer theorem.

Now rescale \(T\) so that its largest eigenspace is its fixed
space.  If \(T\) commuted with all of \(E\), it would be scalar
by irreducibility.  Choose \(h\in E\) with \(c=[T,h]\ne I\).
The rank inequality in \cite[Lemma~3.1, p.~139]{BI:DHP}, or the
identity \([T,h]=T(hT^{-1}h^{-1})\), gives
\[
             \operatorname{rank}(c-I)\le2\nu(T).
 \tag{FTQ.5}
\]
If \(c\) is nonscalar, (FTQ.4) bounds its actual moving rank
below by \(a/40\).  If it is a nonidentity scalar, that rank is
\(a\).  Either case proves \(\nu(T)\ge a/80\).

For the second assertion, put \(S=E/Z(E)\simeq A_m\).
Conjugation defines an injection of the projective normalizer into
\(\operatorname{Aut}(S)\).  Indeed an automorphism of \(E\) which
is trivial on \(S\) has the form \(x\mapsto xz_x\), with
\(z_x\in Z(E)\); centrality makes \(x\mapsto z_x\) a
homomorphism, and perfectness makes it trivial.  An operator
centralizing \(E\) is scalar.  The finite group generated by the
projective images of \(E\) and \(T\) is consequently almost simple
with socle \(A_m\).

The exact input is \cite[Lemma~6.1, pp.~547--548]{BI:GS}:
the parameter \(\alpha(x)\) defined in its Introduction is the
least number of socle-conjugates of \(x\) generating
\(\langle A_m,x\rangle\), and is at most \(m-1\).
Use this statement only at \(m\ge7\).  Rescale \(T\) by an
eigenvalue of maximal multiplicity and choose the corresponding
\(q\le m-1\) conjugates by elements of \(E\).  Let \(L\) be
the group generated by these rescaled lifts.  Its projective image
contains \(S\), so \(E\subset\mathbf C^*L\).  Perfectness gives
\[
            E=[E,E]\subset[\mathbf C^*L,\mathbf C^*L]
                           =[L,L]\subset L.
 \tag{FTQ.6}
\]
The common fixed subspace of the selected operators is therefore
zero: it is fixed by the nontrivial irreducible group \(E\).
Each selected operator has fixed codimension \(\nu(T)\), so
\(a\le q\nu(T)\le(m-1)\nu(T)\).  This proves (FTQ.2).
In particular no unproved lift from projective generation to a
linear perfect core is hidden in the dimension argument.

Finally the cover list in \cite[Chapter~4 opening, p.~31]{BI:HHM}
allows only the double cover for \(A_5\) and divisors of the
sixfold cover for \(A_6\).  Thus \(|E|\le6|A_6|=2160\), and
character orthogonality gives \(a^2\le|E|\), hence \(a\le46\).
Every nonscalar semisimple operator has integral support at least
one, giving the third assertion.  This small-degree argument avoids
any extra interpretation of the exceptional automorphisms at
\(m=6\).
\end{proof}

\subsection*{HD-FT21. Arbitrary numbers and degrees of covered partner pieces}

Suppose the entire partner has an actual factorization
\[
 U=\bigotimes_{j=1}^s U_j,\qquad \dim U_j\ge2,
 \tag{FTQ.7}
\]
whose embedded full matrix algebras are permuted by its actual
finite projective image.  Whenever an operator fixes a factor
algebra, its local operator is required to normalize the indicated
projective core, if one is indicated.  Permit the following pieces:
\begin{enumerate}
\item arbitrary factors of degree at most a specified \(A\ge2\);
\item irreducible modules of represented quasisimple cores, of any
degree and any simple quotient;
\item irreducible modules of a represented central product of
quasisimple cores, with the whole central product normalized;
\item irreducible projective Heisenberg factors as in D-TD08,
with nondegenerate pairing, of any prime and any exponent.
\end{enumerate}
There is no additional residual multiplicity space outside (FTQ.7).
Let \(m_*\) be the largest natural degree of an alternating simple
quotient of any quasisimple component among these pieces, with
\(m_*=1\) if there is none, and define
\[
                         H=\max\{80,A,m_*-1\}.
 \tag{FTQ.8}
\]

\begin{theorem}[HD-FT21]
Every nonscalar operator in this entire partner image has
\[
                         \frac{\nu(T)}{\dim U}\ge\frac1H.
 \tag{FTQ.9}
\]
The degrees and number of quasisimple or Heisenberg pieces are
unrestricted.  The normalization and completeness hypotheses on
(FTQ.7) are part of the assertion.
\end{theorem}

\begin{proof}
The local floors for the first two kinds are \(1/A\) and those
of HD-FT20.  D-TD08 gives \((p-1)/(2p)\ge1/4\) for the fourth
kind, including every compatible nonscalar normalizer and \(p=2\).
This is the separately proved intersection-of-eigenspaces theorem,
not an assumption that a projective normalizer extension splits.

For the third kind, pull the irreducible module back to the direct
product of its quasisimple components.  It is a tensor product
of irreducible modules for those components.  Each is nontrivial,
since the represented central product acts faithfully; each has
degree at least two.  Its factor matrix algebra is generated by
that component.  An automorphism of the central product permutes
these components: modulo the scalar center it permutes the simple
direct factors, and the derived subgroup of the inverse image of
one simple factor recovers its perfect component.  Consequently
every normalizer permutes their full matrix algebras.  HD-FT20
applies to a fixed component.  On a moved component cycle the
weighted-word proof of HD-FT18 gives the quarter floor.  Therefore
the entire irreducible central-product piece has floor at least
\(1/H\).  There is no leftover multiplicity factor in this
irreducible restriction.

Now apply HD-FT18 to all of (FTQ.7).  If an operator fixes every
factor algebra, a nonscalar local factor and spectral convolution
give at least the minimum local floor.  If it moves some factor
algebra, a moved cycle gives the quarter floor independently of
its twists.  Since \(H\ge80\), both alternatives imply (FTQ.9).
\end{proof}

\subsection*{HD-FT22. The tail-two source forces a quadratic jump in an alternating partner}

Retain exactly the actual stable tail-two source of HD-FT13--15:
there is a connected sphere cover with a faithful coefficient image,
an actual quotient onto \(A_n\) or \(S_n\), and
\[
 V=W\otimes U\quad\text{up to scalars},\qquad
 W=S^{(n-2,2)}\ \text{or}\ S^{(n-2,1,1)},\qquad
 r\ge12,\qquad
 \sum_i\operatorname{rank}(1-T_i)=2\dim V+6.
 \tag{FTQ.10}
\]
Transposed labels and ordinary irreducible alternating restrictions
are permitted.  Every original inertia is genuine.  The entire
partner satisfies HD-FT21.  As in HD-FT19, use \(q=4\) when
\(n\ge10^6\), by accepted HD-GC07, or \(q=7\) under the separately
scoped qualitative hypothesis \(n>N_{14}\), \(n\ge64\).
Those results give at least \(q\) genuine identity-projection
entries, whose total normalized moving rank is at most \(61/n^2\).

\begin{theorem}[HD-FT22]
Every such actual source satisfies
\[
                   qn^2\le61\max\{80,A,m_*-1\}.
 \tag{FTQ.11}
\]
In particular, for \(n\ge10^6\), either
\[
             A\ge\frac{4n^2}{61}
       \quad\text{or}\quad
             m_*\ge1+\frac{4n^2}{61}.
 \tag{FTQ.12}
\]
If there are no arbitrary bounded-degree pieces, one may put
\(A=2\).  Thus a source whose distinguished stable tail-two factor
has maximal alternating natural degree is excluded when its entire
partner is built from the covered quasisimple and Heisenberg pieces,
regardless of their dimensions or count.  With bounded-degree pieces
also present and \(m_*\le n\), their bound must satisfy
\(A\ge4n^2/61\).
\end{theorem}

\begin{proof}
At an identity-projection entry the operator is a scalar times
\(I_W\otimes S_i\).  If \(S_i\) is nonscalar, (FTQ.9) bounds its
actual moving rank, including that scalar phase, below by
\(\dim V/H\).  If it is scalar, genuineness of \(T_i\) gives
full moving rank, which is an even stronger bound.  The at least
\(q\) entries and their accepted total budget therefore give
\(q/H\le61/n^2\), proving (FTQ.11).
For \(n\ge10^6\), \(4n^2>61\cdot80\), so the maximum in
(FTQ.11) must be supplied by one of the other two terms; this is
(FTQ.12).  Also \(4n^2>61(n-1)\), giving the maximal-natural-degree
conclusions.  If \(\dim U=1\), all these entries have full rank,
so \(q\le61/n^2\) is already impossible; the empty partner is
covered separately without a fictitious local factor.
\end{proof}

\paragraph{Strict scope enlargement and retained boundary.}
HD-FT19 handled arbitrary standard and ordinary tail-two partner
factors of bounded alternating natural degree.  Here arbitrary
Specht constituents, split constituents, spin modules and represented
cover kernels are allowed.  Nonalternating quasisimple and projective
Heisenberg pieces have no degree or count bound.  Their products
can have arbitrarily large dimension, so this is stronger than the
dimension-only condition \(\dim U\ge4n^2/61\).  No cover or inertia
tuple is inferred merely by selecting such dimensions.

The statement still requires the actual stable distinguished
quotient and the whole partner's stated factor-algebra normalization.
It does not select such a distinguished quotient from a group which
permutes it with other factors, discard an unclassified residual
partner, or claim a complete classification of primitive groups.
All source consequences are necessary conditions; no coefficient-field
descent, polarization, or existence statement is added.
